% Compile with XeLaTeX.
% Protocol 1: faithful transcription of supplied photographs, pp. 162--164.
% Paragraph identifiers continue the existing notes. Protocol 2 review notes added.
\documentclass[11pt,a4paper,twoside]{article}
\usepackage{fontspec}
\IfFontExistsTF{Times New Roman}{\setmainfont{Times New Roman}}{\setmainfont{Latin Modern Roman}}
\usepackage{amsmath,amssymb,mathtools}
\usepackage[a4paper,inner=2.5cm,outer=2cm,top=2.5cm,bottom=2.5cm]{geometry}
\usepackage{setspace}
\setstretch{1.25}
\setlength{\parindent}{2em}
\setlength{\parskip}{5pt}
\setlength{\abovedisplayskip}{6pt}
\setlength{\belowdisplayskip}{6pt}
\usepackage{fancyhdr}
\pagestyle{fancy}
\fancyhf{}
\renewcommand{\headrulewidth}{0pt}
\fancyfoot[RO,LE]{\small $\boldsymbol{\cdot}\;\thepage\;\boldsymbol{\cdot}$}
\newcommand{\pnum}[1]{\makebox[0pt][r]{\scriptsize\textbf{[#1]}\hspace{0.7em}}}
% Traditional Chinese review notes; compile with XeLaTeX.
\usepackage{xeCJK}
\IfFontExistsTF{DFKai-SB}{\setCJKmainfont{DFKai-SB}}{%
  \IfFontExistsTF{Noto Serif CJK TC}{\setCJKmainfont{Noto Serif CJK TC}}{%
    \IfFontExistsTF{Noto Serif TC}{\setCJKmainfont{Noto Serif TC}}{%
      \setCJKmainfont{FandolSong-Regular}}}}
\usepackage[most]{tcolorbox}
\newtcolorbox{studyNote}[1]{enhanced,breakable,
  colback=gray!4,colframe=black!35,boxrule=0.5pt,arc=1mm,
  left=2mm,right=2mm,top=1.5mm,bottom=1.5mm,
  before skip=8pt,after skip=8pt,fonttitle=\bfseries,title={#1}}
\newcommand{\StudyNote}[2]{\begin{studyNote}{#1}#2\end{studyNote}}

\begin{document}
\noindent{\large\bfseries 7.1.3\quad The Second Variation}
\par\vspace{0.6em}\hrule\vspace{0.8em}

% Original page 162; the paragraph continues on page 163.
\noindent\pnum{43}Before we apply the calculus of variations to problems in classical mechanics, there is one thread that we left hanging. We know that the Euler--Lagrange equation is the requirement for a function to be an extremum. But how do we tell if it's a maximum or minimum or something else? It turns out that
% Original page 163.
this information won't actually be needed for our applications to classical mechanics, but we include it here for completeness.

\noindent\pnum{44}To set the scene, we again return to a function of several variables $f(y_1,\ldots,y_n)$. Taylor expanding about a given point $\mathbf{y}$, we have
\begin{equation}
    f(\mathbf{y}+\delta\mathbf{y})-f(\mathbf{y})
    =\frac{\partial f}{\partial y_i}(\mathbf{y})\,\delta y_i
    +\frac12\frac{\partial^2 f}{\partial y_i\partial y_j}(\mathbf{y})\,\delta y_i\delta y_j+\ldots\,.
    \tag{7.28}
\end{equation}

\noindent\pnum{45}We assume that $\mathbf{y}$ is a stationary point, so $\partial f/\partial y_i(\mathbf{y})=0$. That leaves us with the second order variation which is determined by the \textit{Hessian} matrix
\begin{equation}
    H_{ij}(\mathbf{y})=\frac{\partial^2 f}{\partial y_i\partial y_j}(\mathbf{y}).
    \tag{7.29}
\end{equation}

% Protocol 2: stationary point and Hessian.
\StudyNote{複習註解一：stationary point（駐點）}{
    駐點表示「一階變化為零」。對一般函數，條件為 $f'(x_0)=0$；對多變數函數，則為 $\nabla f=0$。
    在變分法中，駐點是一條使泛函的一階變分對所有容許擾動皆為零的路徑：
    \[
        \delta F[y]=0.
    \]
    這\textbf{不代表路徑的斜率 $y'=0$}，也不保證是極大值或極小值；還需要用第二變分判斷。
}
\StudyNote{複習註解二：Hessian 矩陣如何描述二階變化？}{
    Hessian 矩陣把所有二階偏導數排列在一起，扮演一變數函數中 $f''$ 的角色。以兩變數為例，
    \[
        H=\begin{pmatrix}
            \dfrac{\partial^2 f}{\partial y_1^2}           & \dfrac{\partial^2 f}{\partial y_1\partial y_2} \\[2mm]
            \dfrac{\partial^2 f}{\partial y_2\partial y_1} & \dfrac{\partial^2 f}{\partial y_2^2}
        \end{pmatrix}.
    \]
    對角元素描述各座標方向的二階彎曲；混合偏導數描述兩個變數一起改變時的交互影響。若二階偏導數連續，混合偏導數相等，故 $H^T=H$。
    在駐點 $\mathbf y_0$，一階項消失，泰勒展開為
    \[
        f(\mathbf y_0+\delta\mathbf y)-f(\mathbf y_0)
        =\frac12\delta\mathbf y^T H(\mathbf y_0)\delta\mathbf y
        +o(\|\delta\mathbf y\|^2).
    \]
    其中
    \[
        \delta\mathbf y^TH\delta\mathbf y
        =H_{11}(\delta y_1)^2+2H_{12}\delta y_1\delta y_2+H_{22}(\delta y_2)^2.
    \]
    因此，判斷駐點要看\textbf{所有方向}的二階變化，不能只看對角元素。特徵值全正時為嚴格局部極小值；全負時為嚴格局部極大值；正負並存時為鞍點。若含零特徵值且沒有正負並存，二階資訊不足，需進一步分析。
}


\noindent\pnum{46}Because this matrix is symmetric, it can always be diagonalised. This means that there exists a rotation matrix $R\in O(n)$ such that $H'=R^T H R$ is diagonal, with the elements given by the eigenvalues $\lambda_i$ of $H$
\begin{equation}
    H'=\operatorname{diag}(\lambda_1,\ldots,\lambda_n).
    \tag{7.30}
\end{equation}

\noindent\pnum{47}The nature of the stationary point is determined by the signs of these eigenvalues:
\begin{itemize}
    \item If $\lambda_i>0$ for all $i=1,\ldots,n$, then the stationary point is a minimum.
    \item If $\lambda_i<0$ for all $i=1,\ldots,n$, then the stationary point is a maximum.
    \item If $\lambda_i$ are a collection of positive and negative numbers then the stationary point is a saddle point.
    \item If some $\lambda_i=0$ then we have to work harder and look to higher orders to understand what's going on in those directions.
\end{itemize}

% Protocol 2: two worked examples.
\StudyNote{複習註解三：Hessian 的兩個例子}{
    \textbf{例一：$f(y_1,y_2)=y_1^2+3y_2^2$。}
    由 $\nabla f=(2y_1,6y_2)^T=0$，駐點為 $(0,0)$，且
    \[
        H=\begin{pmatrix}2&0\\0&6\end{pmatrix},
        \qquad
        \Delta f=\frac12\delta\mathbf y^TH\delta\mathbf y
        =(\delta y_1)^2+3(\delta y_2)^2.
    \]
    任何非零位移都讓函數值增加，所以原點是嚴格局部極小值，本例也是全域極小值。特徵值 $2,6$ 皆為正；沿 $y_2$ 方向的二階彎曲較強。

    \textbf{例二：$f(y_1,y_2)=y_1^2+4y_1y_2+y_2^2$。}
    由 $2y_1+4y_2=0$ 與 $4y_1+2y_2=0$，駐點也是 $(0,0)$，且
    \[
        H=\begin{pmatrix}2&4\\4&2\end{pmatrix},
        \qquad
        \Delta f=(\delta y_1)^2+4\delta y_1\delta y_2+(\delta y_2)^2.
    \]
    雖然兩個對角元素皆為正，但沿不同方向移動會得到不同符號：
    \[
        \begin{aligned}
            \delta\mathbf y=(t,t)^T  & \quad\Longrightarrow\quad\Delta f=6t^2>0, \\
            \delta\mathbf y=(t,-t)^T & \quad\Longrightarrow\quad\Delta f=-2t^2<0
        \end{aligned}
        \qquad(t\neq0).
    \]
    因此原點是鞍點。由
    \[
        \det(H-\lambda I)=(2-\lambda)^2-16=0
    \]
    可得特徵值 $6,-2$，對應特徵方向分別為 $(1,1)^T$ 與 $(1,-1)^T$。

    兩例皆為二次式，因此 $\Delta f=\tfrac12\delta\mathbf y^TH\delta\mathbf y$ 是精確等式，沒有更高階項。
}


\noindent\pnum{48}Now we extend this to functionals $F[y(x)]$. We vary the path from $y(x)$ to $y(x)+\delta y(x)$ and write
\begin{equation}
    \delta y(x)=\epsilon\eta(x)
    \tag{7.31}
\end{equation}
with $\epsilon\ll1$. We then expand the functional (7.7) to second order
\begin{equation}
    \begin{aligned}
        \delta F[y(x)]={} & \epsilon\left\{
        \int_a^b dx\left[\frac{\partial L}{\partial y}
            -\frac{d}{dx}\left(\frac{\partial L}{\partial y'}\right)\right]\eta(x)
        +\left[\frac{\partial L}{\partial y'}\eta\right]_a^b\right\} \\
                          & +\frac12\epsilon^2\int_a^b dx\left\{
        \frac{\partial^2 L}{\partial y^2}\eta^2
        +2\frac{\partial^2 L}{\partial y\partial y'}\eta\eta'
        +\frac{\partial^2 L}{\partial y'^2}\eta'^2\right\}.
    \end{aligned}
    \tag{7.32}
\end{equation}

\noindent\pnum{49}The first line is what we had before in (7.10). The second line is new. We consider paths between $x=a$ and $x=b$ with fixed end points, so that
% Original page 164: continuation of paragraph [49].
$\eta(a)=\eta(b)=0$. Moreover, we'll take a path that solves the Euler--Lagrange equation. This means that the first line vanishes and we're left just with the second-order term.

\noindent\pnum{50}We have three terms, proportional to $\eta^2$, $\eta\eta'$, and $\eta'^2$. Unfortunately, we don't get to integrate by parts to make all of these look the same. We can, however, use the fact that $2\eta\eta'=(\eta^2)'$ to integrate the second by parts. For an extremum of the action, the second variation is then
\begin{equation}
    \delta^2 F[y(x)]=\frac12\epsilon^2\int_a^b dx\left\{
    \left[\frac{\partial^2 L}{\partial y^2}
        -\frac{d}{dx}\left(\frac{\partial^2 L}{\partial y\partial y'}\right)\right]\eta^2
    +\frac{\partial^2 L}{\partial y'^2}\eta'^2\right\}.
    \tag{7.33}
\end{equation}

\noindent\pnum{51}Here we've thrown away the boundary term on the grounds that $\eta(a)=\eta(b)=0$. We've also renamed this term $\delta^2 F$ to emphasise that it is the second-order variation.

\noindent\pnum{52}The variation (7.33) depends both on the path $y(x)$ that solves the Euler--Lagrange equation, and also on the variation $\eta(x)$. At this point, the story is a more complicated version of what we saw for a function of several variables. Suppose, for example, that, when evaluated on our extremum path $y(x)$, the second-order variation (7.33) is strictly positive:
\begin{equation}
    \delta^2 F[y(x)]>0\quad\text{for all }\eta(x).
    \tag{7.34}
\end{equation}

\noindent\pnum{53}In this case, our extremum was a minimum of the functional. However, it may not be so easy to check that (7.34) holds because we have to do it for all functions $\eta(x)$. There is, however, a simple test that we can perform. If the path $y(x)$ is a local minimum, then it necessarily obeys
\begin{equation}
    \frac{\partial^2 L}{\partial y'^2}\geq0\quad\text{for all }x.
    \tag{7.35}
\end{equation}

\noindent\pnum{54}This is known as the \textit{Legendre condition}. We recognise $\partial^2 L/\partial y'^2$ as the final term in (7.33). The idea behind the Legendre condition is that if (7.35) does \textit{not} hold, so that $\partial^2 L/\partial y'^2<0$ for some region, then we can always find a test function so that $\eta$ is small in that region while $\eta'$ is large and make the final term in (7.33) arbitrarily negative so that $\delta^2 F<0$. The Legendre condition is therefore necessary for the path to be a minimum. It is not, however, sufficient.

% End of section 7.1.3. No following-section text is included.
\end{document}
